\documentclass[../main.tex]{subfiles}
\begin{document}

\chapter{Optimization}


\section{Stochastic optimization}

We consider optimization of stochastic objectives of the form
\[
\mathcal{L}(\boldsymbol{\theta}) = \mathbb{E}_{q_{\boldsymbol{\theta}}(\boldsymbol{z})}\left[ \tilde{L}(\boldsymbol{\theta}, \boldsymbol{z}) \right] 
\]
where $\boldsymbol{z}$ is a r.v., such s an external noise


\subsection{SGD for optimizing the parameters of a distribution}

$\boldsymbol{z}$ could be a latent variable sampled from an inference network $q_{\boldsymbol{\theta}}$.
\[
\begin{aligned}
\nabla_{\boldsymbol{\theta}} \mathbb{E}_{q_{\boldsymbol{\theta}}(\boldsymbol{z})}\left[ \tilde{L}(\boldsymbol{\theta}, \boldsymbol{z}) \right]  &= \nabla_{\boldsymbol{\theta}} \int \mathcal{L}(\boldsymbol{\theta}, \boldsymbol{z}) q_{\boldsymbol{\theta}}(\boldsymbol{z}) d\boldsymbol{z} = \int \nabla _{\boldsymbol{\theta}} \tilde{L}(\boldsymbol{\theta}, \boldsymbol{z}) q_{\boldsymbol{\theta}}(\boldsymbol{z}) d \boldsymbol{z} \\
&= \int \left[ \nabla_{\boldsymbol{\theta}} \tilde{\mathcal{L}}(\boldsymbol{\theta}, \boldsymbol{z}) \right] q_{\boldsymbol{\theta}}(\boldsymbol{z}) d \boldsymbol{z} + \int \tilde{L}(\boldsymbol{\theta}, \boldsymbol{z}) \left[ \nabla_{\boldsymbol{\theta}} q_{\boldsymbol{\theta}}(\boldsymbol{z}) \right] d \boldsymbol{z}
\end{aligned}
\]

The first term can be apprximated by
\[
\int \left[ \nabla_{\boldsymbol{\theta}} \tilde{\mathcal{L}} (\boldsymbol{\theta}, \boldsymbol{z})\right] q_{\boldsymbol{\theta}}(\boldsymbol{z}) d \boldsymbol{z} \approx  \frac{1}{S} \sum_{s=1}^{S} \nabla_{\boldsymbol{\theta}} \tilde{L}(\boldsymbol{\theta}, z_{s}).
\]
where $\boldsymbol{z}_{s} \sim  q_{\boldsymbol{\theta}}$. 

Consider the second term, 
\[
I  \overset{\underset{\mathrm{def}}{}}{=} \int \tilde{\mathcal{L}}(\boldsymbol{\theta}, \boldsymbol{z}) \left[ \nabla_{\boldsymbol{\theta}} q_{\boldsymbol{\theta}}(\boldsymbol{z}) \right] d\boldsymbol{z}
\]
There are two main methods to approximate $I$.

\subsection{Score function estimator (REINFORCE)}

\[
\nabla_{\boldsymbol{\theta}} q_{\boldsymbol{\theta}}(\boldsymbol{z}) = q_{\boldsymbol{\theta}}(\boldsymbol{z}) \nabla_{\boldsymbol{\theta}} \log q_{\boldsymbol{\theta}}(\boldsymbol{z})
\]
So
\[
I = \mathbb{E}_{q_{\boldsymbol{\theta}}(\boldsymbol{z})} \left[ \tilde{\mathcal{L}}(\boldsymbol{\theta}, \boldsymbol{z}) \nabla_{\boldsymbol{\theta}} \log q_{\boldsymbol{\theta}}(\boldsymbol{z}) \right] .
\]
which can be approxiated by Monte Carlo.


\subsection{Reparameterization trick}\label{sec:reparameterization}

Score function estimator can have \emph{high variance}. 
\assumption{}{} {
   \begin{itemize}
   \item  $\tilde{\mathcal{L}}$ is differentiable wrt $\boldsymbol{z}$
   \item We can compute a sample from $q_{\boldsymbol{\theta}}(\boldsymbol{z})$ by first sampling $\boldsymbol{\epsilon}$ from some noise distribution $q_{0}$
   which is independent of $\boldsymbol{\theta}$, and then transforming $\boldsymbol{z}$ using a deterministic and differentiable function $\boldsymbol{z} = g(\boldsymbol{\theta}, \boldsymbol{\epsilon})$.
   \end{itemize} 
}

\[
\mathcal{L}(\boldsymbol{\theta}) = \mathbb{E}_{q_{0}(\boldsymbol{\epsilon})} \left[ \tilde{\mathcal{L}}(\boldsymbol{\theta}, g(\boldsymbol{\theta}, \boldsymbol{\epsilon})) \right] 
\]

\[
\nabla_{\boldsymbol{\theta}} \mathcal{L}(\boldsymbol{\theta}) = \mathbb{E}_{q_0(\boldsymbol{\epsilon})} \left[ \nabla_{\boldsymbol{\theta}} \tilde{\mathcal{L}} (\boldsymbol{\theta}, g(\boldsymbol{\theta}, \boldsymbol{\epsilon}))\right] \approx  \frac{1}{S} \sum_{s=1}^{S} \nabla_{\boldsymbol{\theta}} \tilde{\mathcal{L}} \left( \boldsymbol{\theta}, g(\boldsymbol{\theta}, \boldsymbol{\epsilon}_{s}) \right) .
\]
where $\boldsymbol{\epsilon}_{s} \sim  q_{0}$.

\begin{note}[A example]
   Let $z \sim i \mathcal{N}(\mu, v)$, $v = \sigma^{2}$.
   \[
   \epsilon \sim  \mathcal{N}(0,1), \quad z = \mu + \sqrt{v} \epsilon
   \]

   \[
   \mathcal{L}(\mu, v) = \mathbb{E}_{z \sim  \mathcal{N}(\mu, v)} \left[ \tilde{\mathcal{L}}(z) \right]  = \mathbb{E}_{\epsilon \sim  \mathcal{N}(0,1)} \left[ \tilde{\mathcal{L}} (\mu + \sqrt{v} \epsilon) \right] 
   \]
   Now the stochasticity is independent of $\mu$ and $v$, so we can compute the gradient by backpropagation.
   Let $z(\mu, v, \epsilon) = \mu + \sqrt{v} \epsilon$, 
   \[
   \frac{\partial \mathcal{L}}{\partial \mu} = \frac{\partial }{\partial \mu} \mathbb{E}_{\epsilon} \left[ \tilde{L}(z) \right] = \mathbb{E}_{\epsilon} \left[ \frac{\partial \tilde{\mathcal{L}}}{\partial z} \cdot  \frac{\partial z}{\partial \mu} \right]  = \mathbb{E}_{\epsilon} \left[ \frac{\partial \tilde{\mathcal{L}}(z)}{\partial z} \right] 
   \]

   \[
   \frac{\partial \mathcal{L}}{\partial v} = \mathbb{E}_{\epsilon} \left[ \frac{\partial \tilde{L}}{\partial z} \cdot  \frac{\partial z}{\partial v} \right]  = \mathbb{E}_{\epsilon} \left[ \frac{\partial \tilde{L}(z)}{\partial z} \cdot  \frac{\epsilon}{2 \sqrt{ v}} \right] 
   \]
   Then we can use Monte Carlo to compute this.
\end{note} 








\end{document}
